\documentclass[10pt,a4paper]{article}
\usepackage[utf8]{inputenc}
\usepackage[T1]{fontenc}
\usepackage[spanish]{babel}
\usepackage{amsmath,amssymb}
\usepackage{geometry}
\geometry{top=1.8cm,bottom=1.8cm,left=2cm,right=2cm}
\usepackage{booktabs}
\usepackage{siunitx}
\sisetup{output-decimal-marker={,}}
\usepackage{multicol}
\usepackage[hidelinks]{hyperref}
\setlength{\parindent}{0pt}
\setlength{\parskip}{4pt}

\begin{document}

\begin{center}
  {\large\textbf{Formulario — Mecánica y Ondas II (Laboratorio)}}\\[2pt]
  {\small Luis López Nasser · Universidad UNIE · Curso 2025--26}
\end{center}

\hrule\vspace{6pt}

\begin{multicols}{2}

% ─────────────────────────────────────────────────────────
\section*{Práctica: Ondas estacionarias en una cuerda}
% ─────────────────────────────────────────────────────────

\subsection*{Ecuación de ondas}
\begin{equation*}
  \frac{\partial^2 y}{\partial x^2} = \frac{1}{v^2}\frac{\partial^2 y}{\partial t^2}
\end{equation*}
Válida para perturbaciones de pequeña amplitud ($\alpha \ll 1$).

\subsection*{Velocidad de fase}
\begin{equation*}
  v = \sqrt{\frac{T}{\mu}}, \qquad \mu = \frac{m_c}{L_T}
\end{equation*}
$T$: tensión (N); $\mu$: densidad lineal (\si{\kg\per\m}).

\subsection*{Onda estacionaria}
\begin{equation*}
  y(x,t) = 2A\sin(kx)\sin(\omega t)
\end{equation*}
Superposición de dos viajeras opuestas con desfase $\pi$.

\subsection*{Modos normales (extremos fijos)}
\begin{equation*}
  \boxed{f_n = \frac{n\,v}{2L} = \frac{n}{2L}\sqrt{\frac{T}{\mu}}}, \quad n = 1,2,3,\ldots
\end{equation*}
Modo $n$: $n$ vientres, $n-1$ nodos interiores.

\subsection*{Velocidad directa (un modo)}
\begin{equation*}
  v = \frac{2L\,f_n}{n}
\end{equation*}

\subsection*{Velocidad por ajuste lineal}
Ajuste $f_n = an + b$ $\Rightarrow$ $v_a = 2La$.
\begin{equation*}
  \delta v_a = 2L\,\delta a
\end{equation*}
donde $\delta a$ es el error estándar de la pendiente del ajuste.

\subsection*{Propagación de incertidumbre en $v$}
\begin{equation*}
  \delta v = \frac{2}{n}\sqrt{(f_n\,\delta L)^2 + (L\,\delta f)^2}
\end{equation*}
El término dominante es $(L\,\delta f)^2$.

\subsection*{Tensión aplicada}
\begin{equation*}
  T_i = (m_s + 0{,}010\cdot i)\,g, \quad \delta T = g\,\delta m
\end{equation*}
$m_s$: masa del soporte; $g = \SI{9.81}{\m\per\s\squared}$.

\subsection*{Efecto de doblar la longitud}
\begin{equation*}
  L \to 2L \implies f_n' = \frac{f_n}{2}
\end{equation*}
$v$ no cambia; $\lambda_n$ se duplica; $f_n$ se divide entre dos.

\subsection*{Datos del experimento}
\begin{center}
  \small
  \begin{tabular}{ll}
    \toprule
    Longitud útil $L$   & \SI{0.1315}{\m} \\
    Longitud total $L_T$& \SI{0.82}{\m}   \\
    $\delta L$          & \SI{0.5}{\mm}   \\
    $\delta f$          & \SI{0.5}{\Hz}   \\
    Modos medidos       & $n = 3$ a $8$   \\
    \bottomrule
  \end{tabular}
\end{center}

\subsection*{Resultados principales}
\begin{center}
  \small
  \begin{tabular}{lcc}
    \toprule
    Tensión & $\bar{v}$ (m/s) & $v_a \pm \delta v_a$ (m/s) \\
    \midrule
    $T_0 = \SI{0.098}{\N}$ & 0,791 & $0{,}885\pm0{,}059$ \\
    $T_1 = \SI{0.196}{\N}$ & 0,927 & $1{,}105\pm0{,}102$ \\
    $T_2 = \SI{0.294}{\N}$ & 1,149 & $1{,}207\pm0{,}039$ \\
    $T_3 = \SI{0.392}{\N}$ & 1,264 & $1{,}268\pm0{,}018$ \\
    \bottomrule
  \end{tabular}
\end{center}

\end{multicols}

\end{document}
